\clearpage
\VerdierPhysicalPageAnchor{0042}%
\VerdierSemanticAnchor{ch1-1-3-7-cont}
只使用 $G^\circ$-范畴的语言。不过，我们不会完全忘掉“分次范畴”的观点，
并采用下列术语：

\VerdierSemanticAnchor{ch1-1-3-8}
\noindent\textbf{定义 1.3.8.}\quad
设 $G$ 为一个群，$(\mathcal C,G^\circ,T,c)$ 为一个 $G^\circ$-范畴。设
$g,h$ 为 $G$ 的两个元素，$X,Y$ 为 $\mathcal C$ 的两个对象。$\mathcal C$ 中的
一个态射

\VerdierSemanticAnchor{ch1-1-3-8-morphism}
\[
  u\colon T(h)X\longrightarrow T(gh)Y
\]
称为\emph{从 $X$ 到 $Y$ 的次数为 $g$ 的态射}。

下面说明这个术语不会引起混淆。事实上，对 $G$ 中每一对元素 $h,h'$，函子
$T(h^{-1}h')$ 以及同构 $c(h,h^{-1}h')$、$c(gh,h^{-1}h')$ 定义一个双函子的
同构

\VerdierSemanticAnchor{ch1-1-3-8-eq1}
\begin{equation}
  i_g(h,h')\colon
  \operatorname{Hom}_{\mathcal C}\bigl(T(h)X,T(gh)Y\bigr)
  \overset{\sim}{\longrightarrow}
  \operatorname{Hom}_{\mathcal C}\bigl(T(h')X,T(gh')Y\bigr).
  \tag{1.3.8.1}
\end{equation}

余循环恒等式 (1.2.1.3) 表明，对每组三个元素 $h,h',h''\in G$，都有关系式

\VerdierSemanticAnchor{ch1-1-3-8-eq2}
\begin{equation}
  i_g(h,h'')=i_g(h',h'')i_g(h,h'),
  \tag{1.3.8.2}
\end{equation}
因此，如果借助同构 $i_g(h,h')$ 来等同各个集合

\VerdierSemanticAnchor{ch1-1-3-8-identification}
\[
  \operatorname{Hom}_{\mathcal C}\bigl(T(h)X,T(gh)Y\bigr),
\]
那么公式 (1.3.8.2) 表明这些不同的等同彼此相容。此外，对每四个元素
$g,g',h,h'\in G$ 以及 $\mathcal C$ 的每三个对象 $X,Y,Z$，下图 (1.3.8.3)
交换：

\VerdierSemanticAnchor{ch1-1-3-8-eq3}
\begin{center}
  \resizebox{0.92\textwidth}{!}{%
    $\begin{tikzcd}[ampersand replacement=\&,column sep=5em,row sep=4em]
      \operatorname{Hom}_{\mathcal C}\bigl(T(h)X,T(gh)Y\bigr)
      \times
      \operatorname{Hom}_{\mathcal C}\bigl(T(gh)Y,T(g'gh)Z\bigr)
      \arrow[r]
      \arrow[d,
        "{i_g(h,h')\times i_{g'}(gh,gh')}"',
        "\sim"]
      \&
      \operatorname{Hom}_{\mathcal C}\bigl(T(h)X,T(g'gh)Z\bigr)
      \arrow[d,"{i_{g'g}(h,h')}"',"\sim"]
      \\
      \operatorname{Hom}_{\mathcal C}\bigl(T(h')X,T(gh')Y\bigr)
      \times
      \operatorname{Hom}_{\mathcal C}\bigl(T(gh')Y,T(g'gh')Z\bigr)
      \arrow[r]
      \&
      \operatorname{Hom}_{\mathcal C}\bigl(T(h')X,T(g'gh')Z\bigr)
    \end{tikzcd}$%
  }
  \hfill\textup{(1.3.8.3)}
\end{center}

\VerdierSemanticAnchor{ch1-1-3-8-composition}
（横箭头由 $\mathcal C$ 中的复合定义。）因此，图 (1.3.8.3) 使我们能够无歧义
地定义两个次数分别为 $g$ 和 $g'$ 的态射的复合；更精确地说，是在同构
$i_g(h,h')$ 所定义的等同之下无歧义地定义。立即可以验证，这个复合定义与
